Srovnání míry \ac{AROPE} podle skupin zemí (Severní Evropa, ZE, Visegrád) zdůrazňuje, že skupinová příslušnost k~modelu trhu práce je silnějším prediktorem \ac{AROPE} než samotná výše \ac{HDP}: severské země s~koordinovanými mzdovými systémy a~silnou \ac{APZ} dosahují hodnot pod \SI{15}{\percent}, zatímco Visegrád osciluje kolem \SIrange{20}{30}{\percent}.
Skupinový přístup dále ukazuje divergenci uvnitř V4: PL a~HU se v~poslední dekádě přibližovaly k~průměru EU27, zatímco CZ vývoj stagnuje na relativně nízké, ale strukturálně zranitelné úrovni.
Klíčovým faktorem rozdílů mezi skupinami není pouze úroveň \ac{HDP}, ale intenzita kolektivního vyjednávání a~výdaje na aktivní politiku zaměstnanosti, které v~severských zemích fungují jako pojistka pro pracovníky vycházející z~flexibilní pracovní smlouvy (flexicurity).
Přiblížení ČR k~severské skupině by vyžadovalo souběžně zvýšit \ac{KV} pokrytí i~\ac{APZ} výdaje --- jedná se tedy o~systémový balíček, nikoli o~izolovanou reformu.

Srovnání míry \ac{AROPE} podle skupin zemí (Severní Evropa, ZE, Visegrád) zdůrazňuje, že skupinová příslušnost k~modelu trhu práce je silnějším prediktorem \ac{AROPE} než samotná výše \ac{HDP}: severské země s~koordinovanými mzdovými systémy a~silnou \ac{APZ} dosahují hodnot pod \SI{15}{\percent}, zatímco Visegrád osciluje kolem \SIrange{20}{30}{\percent}.
Skupinový přístup dále ukazuje divergenci uvnitř V4: PL a~HU se v~poslední dekádě přibližovaly k~průměru EU27, zatímco CZ vývoj stagnuje na relativně nízké, ale strukturálně zranitelné úrovni.
Klíčovým faktorem rozdílů mezi skupinami není pouze úroveň \ac{HDP}, ale intenzita kolektivního vyjednávání a~výdaje na aktivní politiku zaměstnanosti, které v~severských zemích fungují jako pojistka pro pracovníky vycházející z~flexibilní pracovní smlouvy (flexicurity).
Přiblížení ČR k~severské skupině by vyžadovalo souběžně zvýšit \ac{KV} pokrytí i~\ac{APZ} výdaje --- jedná se tedy o~systémový balíček, nikoli o~izolovanou reformu.

Srovnání míry \ac{AROPE} podle skupin zemí (Severní Evropa, ZE, Visegrád) zdůrazňuje, že skupinová příslušnost k~modelu trhu práce je silnějším prediktorem \ac{AROPE} než samotná výše \ac{HDP}: severské země s~koordinovanými mzdovými systémy a~silnou \ac{APZ} dosahují hodnot pod \SI{15}{\percent}, zatímco Visegrád osciluje kolem \SIrange{20}{30}{\percent}.
Skupinový přístup dále ukazuje divergenci uvnitř V4: PL a~HU se v~poslední dekádě přibližovaly k~průměru EU27, zatímco CZ vývoj stagnuje na relativně nízké, ale strukturálně zranitelné úrovni.
Klíčovým faktorem rozdílů mezi skupinami není pouze úroveň \ac{HDP}, ale intenzita kolektivního vyjednávání a~výdaje na aktivní politiku zaměstnanosti, které v~severských zemích fungují jako pojistka pro pracovníky vycházející z~flexibilní pracovní smlouvy (flexicurity).
Přiblížení ČR k~severské skupině by vyžadovalo souběžně zvýšit \ac{KV} pokrytí i~\ac{APZ} výdaje --- jedná se tedy o~systémový balíček, nikoli o~izolovanou reformu.

Srovnání míry \ac{AROPE} podle skupin zemí (Severní Evropa, ZE, Visegrád) zdůrazňuje, že skupinová příslušnost k~modelu trhu práce je silnějším prediktorem \ac{AROPE} než samotná výše \ac{HDP}: severské země s~koordinovanými mzdovými systémy a~silnou \ac{APZ} dosahují hodnot pod \SI{15}{\percent}, zatímco Visegrád osciluje kolem \SIrange{20}{30}{\percent}.
Skupinový přístup dále ukazuje divergenci uvnitř V4: PL a~HU se v~poslední dekádě přibližovaly k~průměru EU27, zatímco CZ vývoj stagnuje na relativně nízké, ale strukturálně zranitelné úrovni.
Klíčovým faktorem rozdílů mezi skupinami není pouze úroveň \ac{HDP}, ale intenzita kolektivního vyjednávání a~výdaje na aktivní politiku zaměstnanosti, které v~severských zemích fungují jako pojistka pro pracovníky vycházející z~flexibilní pracovní smlouvy (flexicurity).
Přiblížení ČR k~severské skupině by vyžadovalo souběžně zvýšit \ac{KV} pokrytí i~\ac{APZ} výdaje --- jedná se tedy o~systémový balíček, nikoli o~izolovanou reformu.

\input{texparts/python/arope_groups}



